%-------------------------------------------------------------------------------
%	SECTION TITLE
%-------------------------------------------------------------------------------
\cvsection{Research Experience}


%-------------------------------------------------------------------------------
%	CONTENT
%-------------------------------------------------------------------------------
\begin{cventries}

%---------------------------------------------------------
  \cventry
    {Undergraduate Research Assistant, Advisor: Prof. Patrick McDaniel} % Job title
    {MadS\&P - Security and Privacy Research Group at UW-Madison} % Organization
    {Madison, WI, U.S.A} % Location
    {Sep. 2024 - Present} % Date(s)
    {
      \begin{cvitems} % Description(s) of tasks/responsibilities
        \item {Conducting scientific measurements and evaluations of the performance advantages of \textbf{LibIHT} over existing dynamic binary tracing framework, with a focus on the runtime overhead, tracing accuracy, and resilience to anti-tracing techniques.}
        \item {Developing a comprehensive evaluation methodology to assess the performance and effectiveness in real-world reverse engineering scenarios, including large-scale softwares and malware samples, to provide insights on practical implications of security research and industry applications.}
        \item {Contributing to a research paper to present the findings and insights from the performance evaluation, targeting top-tier security conferences.}
        % Shortened version by ChatGPT
        % Evaluating performance benefits of LibHT over dynamic binary tracing frameworks, focusing on runtime, accuracy, and anti-tracing resilience.
        % Developing methodologies for performance analysis of large-scale software and malware in reverse engineering scenarios.
        % Contributing to a research paper targeting top-tier security conferences.
      \end{cvitems}
    }

%---------------------------------------------------------
  \cventry
    {Undergraduate Research Assistant, Advisor: Prof. Yan Shoshitaishvili} % Job title
    {SEFCOM - Laboratory of Security Engineering for Future Computing at ASU} % Organization
    {Tempe, AZ, U.S.A} % Location
    {May 2024 - Aug. 2024} % Date(s)
    {
      \begin{cvitems} % Description(s) of tasks/responsibilities
        \item {Assisted in the design and development of \textbf{Oxidizer}, a Rust-specific decompiler on top of \textit{angr} C decompilation pipeline, aimed at producing Rust pseudo-code with recovered control flow structures, function prototypes, and common-used macros and struct data types from Rust binaries.}
        \item {Developed a Rust type scanner in Rust to collect type information from Rust standard libraries and popular crates, and integrated the type information into the decompilation process to recover Rust-specific data types and function signatures with pattern matching and type inference algorithms.}
        \item {Collected a comprehensive dataset of Rust binaries with source code to perform evaluations on the decompilation accuracy and effectiveness.}
        % Shortened version by ChatGPT
        % Developed Oxidizer, a Rust-specific decompiler, enhancing control flow recovery, function prototypes, and Rust-specific macros.
        % Built a Rust type scanner to extract and integrate type information into decompilation, improving data type and function signature recovery.
        % Curated a dataset of Rust binaries for evaluating decompilation accuracy and effectiveness.
      \end{cvitems}
    }

%---------------------------------------------------------
  \cventry
    {Undergraduate Research Assistant, Advisor: Prof. Yan Shoshitaishvili} % Job title
    {SEFCOM - Laboratory of Security Engineering for Future Computing at ASU} % Organization
    {Tempe, AZ, U.S.A} % Location
    {May 2023 - Jul. 2023} % Date(s)
    {
      \begin{cvitems} % Description(s) of tasks/responsibilities
        \item {Engineered \textbf{LibHopper}, an automated tool for discovering exploitation primitives of the internal state structures in user libraries.}
        \item {Designed analyzation pipeline consisting of debugger scripts to setup context for the internal state structures, and symbolic execution engine to evaluate security implications on memory corruption and control flow hijack, reducing manual analysis time and providing structured assessments.}
        \item {Concluded a comprehensive dataset of potential exploitation primitives and possiable mitigation strategies in popular user libraries, including \textit{openssl}, \textit{openssh}, \textit{libcurl}, and \textit{libpng}, etc., contributed to the development of automated exploit generation tools and vulnerability scanners.}
        % Shortened version by ChatGPT
        % Developed LibHopper, a tool for discovering exploitation primitives in user libraries.
        % Created a pipeline using debugger scripts and symbolic execution for analyzing internal structures and evaluating memory corruption and control flow vulnerabilities.
        % Built a dataset of exploitation primitives and mitigation strategies for libraries (e.g., OpenSSL, libcurl), aiding automated exploit and vulnerability tool development.
      \end{cvitems}
    }

%---------------------------------------------------------
\end{cventries}
